\section{Intervalos de confianza e incertidumbre predictiva}

\subsection{La proyección no es certeza}

Una proyección estadística es una estimación condicionada a datos históricos, modelo seleccionado y supuestos metodológicos. No es una certeza ni una promesa de ocurrencia. Esta aclaración es indispensable en análisis de costos de obras civiles porque decisiones contractuales o presupuestales pueden verse afectadas si se interpreta el pronóstico como valor garantizado.

El software reporta un valor puntual proyectado y un intervalo aproximado. El valor puntual corresponde a la estimación central del modelo. El intervalo comunica incertidumbre: un rango razonable dentro del cual podría ubicarse el índice futuro bajo las condiciones históricas observadas y el procedimiento estadístico usado.

\subsection{Error de predicción}

El error de predicción se define como:

\begin{equation}
    \varepsilon_{T+h} = y_{T+h} - \hat{y}_{T+h},
    \label{eq:error_prediccion}
\end{equation}

donde \(h\) es el horizonte de pronóstico. Al momento de proyectar, \(y_{T+h}\) no se conoce. Por tanto, el intervalo busca anticipar la variabilidad esperada de ese error.

En general, la incertidumbre aumenta con el horizonte. Proyectar el mes siguiente suele ser más defendible que proyectar varios años adelante. En índices de construcción, esta incertidumbre incluye no solo variación estadística, sino riesgo económico: cambios de precios internacionales, conflictos logísticos, políticas públicas, tasa de cambio, oferta de materiales y condiciones de demanda.

\subsection{Intervalos aproximados por residuos}

Una aproximación práctica consiste en usar la distribución de residuos históricos para estimar incertidumbre futura. Si los residuos representan errores típicos del modelo, pueden emplearse para construir intervalos empíricos. El software utiliza una aproximación tipo bootstrap residual, en la que se remuestrean residuos para generar una distribución del valor proyectado.

Conceptualmente, si \(\hat{y}_{T+h}\) es la proyección puntual y \(e^{\ast}\) es un residuo remuestreado, una predicción simulada puede expresarse como:

\begin{equation}
    \hat{y}^{\ast}_{T+h} = \hat{y}_{T+h} + e^{\ast}.
    \label{eq:bootstrap_residual}
\end{equation}

Al repetir este procedimiento muchas veces, se obtienen percentiles empíricos. Un intervalo del 95\% puede aproximarse con los percentiles 2.5 y 97.5:

\begin{equation}
    IC_{95\%} = \left[P_{2.5}(\hat{y}^{\ast}_{T+h}), P_{97.5}(\hat{y}^{\ast}_{T+h})\right].
    \label{eq:ic_percentiles}
\end{equation}

Este enfoque es útil porque no exige una normalidad perfecta de residuos. Sin embargo, sigue dependiendo de que los errores históricos sean informativos sobre errores futuros.

\subsection{Interpretación práctica del intervalo}

Si el informe indica:

\begin{center}
\textit{Índice proyectado: 132.4; IC 95\%: [130.9, 134.1],}
\end{center}

la lectura correcta es: bajo el modelo y los datos usados, existe un rango de incertidumbre asociado a la predicción. No significa que haya una probabilidad mecánica del 95\% de que el valor oficial futuro caiga allí en todos los escenarios económicos posibles. Significa que, bajo el procedimiento estadístico adoptado, ese rango resume la variabilidad histórica compatible con el modelo.

\subsection{Amplitud del intervalo}

La amplitud del intervalo se calcula como:

\begin{equation}
    A_{IC} = L_{sup} - L_{inf},
    \label{eq:amplitud_ic}
\end{equation}

donde \(L_{sup}\) y \(L_{inf}\) son los límites superior e inferior. Una amplitud grande puede indicar alta volatilidad, mal ajuste, pocos datos o incertidumbre elevada. En la selección automática de modelos, un intervalo excesivamente amplio reduce la utilidad práctica del pronóstico, aunque el modelo tenga buen ajuste histórico.

\subsection{Riesgo de extrapolación}

Extrapolar significa proyectar fuera del rango observado. Toda proyección futura es una extrapolación temporal. El riesgo aumenta cuando:

\begin{itemize}
    \item el horizonte es largo;
    \item la serie tuvo cambios recientes bruscos;
    \item el modelo seleccionado es muy curvo;
    \item los residuos muestran autocorrelación;
    \item el intervalo de predicción es amplio;
    \item existen factores económicos externos no incorporados.
\end{itemize}

Por ello, el informe técnico debe recomendar que las proyecciones sean usadas como soporte cuantitativo, no como sustituto del análisis económico y contractual.
